% Options for packages loaded elsewhere
% Options for packages loaded elsewhere
\PassOptionsToPackage{unicode}{hyperref}
\PassOptionsToPackage{hyphens}{url}
\PassOptionsToPackage{dvipsnames,svgnames,x11names}{xcolor}
%
\documentclass[
  12pt,
  letterpaper,
  DIV=11,
  numbers=noendperiod]{scrartcl}
\usepackage{xcolor}
\usepackage[left=2.5cm,right=2.5cm,top=2.5cm,bottom=2.5cm]{geometry}
\usepackage{amsmath,amssymb}
\setcounter{secnumdepth}{5}
\usepackage{iftex}
\ifPDFTeX
  \usepackage[T1]{fontenc}
  \usepackage[utf8]{inputenc}
  \usepackage{textcomp} % provide euro and other symbols
\else % if luatex or xetex
  \usepackage{unicode-math} % this also loads fontspec
  \defaultfontfeatures{Scale=MatchLowercase}
  \defaultfontfeatures[\rmfamily]{Ligatures=TeX,Scale=1}
\fi
\usepackage{lmodern}
\ifPDFTeX\else
  % xetex/luatex font selection
  \setmainfont[]{Times New Roman}
\fi
% Use upquote if available, for straight quotes in verbatim environments
\IfFileExists{upquote.sty}{\usepackage{upquote}}{}
\IfFileExists{microtype.sty}{% use microtype if available
  \usepackage[]{microtype}
  \UseMicrotypeSet[protrusion]{basicmath} % disable protrusion for tt fonts
}{}
\makeatletter
\@ifundefined{KOMAClassName}{% if non-KOMA class
  \IfFileExists{parskip.sty}{%
    \usepackage{parskip}
  }{% else
    \setlength{\parindent}{0pt}
    \setlength{\parskip}{6pt plus 2pt minus 1pt}}
}{% if KOMA class
  \KOMAoptions{parskip=half}}
\makeatother
% Make \paragraph and \subparagraph free-standing
\makeatletter
\ifx\paragraph\undefined\else
  \let\oldparagraph\paragraph
  \renewcommand{\paragraph}{
    \@ifstar
      \xxxParagraphStar
      \xxxParagraphNoStar
  }
  \newcommand{\xxxParagraphStar}[1]{\oldparagraph*{#1}\mbox{}}
  \newcommand{\xxxParagraphNoStar}[1]{\oldparagraph{#1}\mbox{}}
\fi
\ifx\subparagraph\undefined\else
  \let\oldsubparagraph\subparagraph
  \renewcommand{\subparagraph}{
    \@ifstar
      \xxxSubParagraphStar
      \xxxSubParagraphNoStar
  }
  \newcommand{\xxxSubParagraphStar}[1]{\oldsubparagraph*{#1}\mbox{}}
  \newcommand{\xxxSubParagraphNoStar}[1]{\oldsubparagraph{#1}\mbox{}}
\fi
\makeatother


\usepackage{longtable,booktabs,array}
\newcounter{none} % for unnumbered tables
\usepackage{calc} % for calculating minipage widths
% Correct order of tables after \paragraph or \subparagraph
\usepackage{etoolbox}
\makeatletter
\patchcmd\longtable{\par}{\if@noskipsec\mbox{}\fi\par}{}{}
\makeatother
% Allow footnotes in longtable head/foot
\IfFileExists{footnotehyper.sty}{\usepackage{footnotehyper}}{\usepackage{footnote}}
\makesavenoteenv{longtable}
\usepackage{graphicx}
\makeatletter
\newsavebox\pandoc@box
\newcommand*\pandocbounded[1]{% scales image to fit in text height/width
  \sbox\pandoc@box{#1}%
  \Gscale@div\@tempa{\textheight}{\dimexpr\ht\pandoc@box+\dp\pandoc@box\relax}%
  \Gscale@div\@tempb{\linewidth}{\wd\pandoc@box}%
  \ifdim\@tempb\p@<\@tempa\p@\let\@tempa\@tempb\fi% select the smaller of both
  \ifdim\@tempa\p@<\p@\scalebox{\@tempa}{\usebox\pandoc@box}%
  \else\usebox{\pandoc@box}%
  \fi%
}
% Set default figure placement to htbp
\def\fps@figure{htbp}
\makeatother


% definitions for citeproc citations
\NewDocumentCommand\citeproctext{}{}
\NewDocumentCommand\citeproc{mm}{%
  \begingroup\def\citeproctext{#2}\cite{#1}\endgroup}
\makeatletter
 % allow citations to break across lines
 \let\@cite@ofmt\@firstofone
 % avoid brackets around text for \cite:
 \def\@biblabel#1{}
 \def\@cite#1#2{{#1\if@tempswa , #2\fi}}
\makeatother
\newlength{\cslhangindent}
\setlength{\cslhangindent}{1.5em}
\newlength{\csllabelwidth}
\setlength{\csllabelwidth}{3em}
\newenvironment{CSLReferences}[2] % #1 hanging-indent, #2 entry-spacing
 {\begin{list}{}{%
  \setlength{\itemindent}{0pt}
  \setlength{\leftmargin}{0pt}
  \setlength{\parsep}{0pt}
  % turn on hanging indent if param 1 is 1
  \ifodd #1
   \setlength{\leftmargin}{\cslhangindent}
   \setlength{\itemindent}{-1\cslhangindent}
  \fi
  % set entry spacing
  \setlength{\itemsep}{#2\baselineskip}}}
 {\end{list}}
\usepackage{calc}
\newcommand{\CSLBlock}[1]{\hfill\break\parbox[t]{\linewidth}{\strut\ignorespaces#1\strut}}
\newcommand{\CSLLeftMargin}[1]{\parbox[t]{\csllabelwidth}{\strut#1\strut}}
\newcommand{\CSLRightInline}[1]{\parbox[t]{\linewidth - \csllabelwidth}{\strut#1\strut}}
\newcommand{\CSLIndent}[1]{\hspace{\cslhangindent}#1}



\setlength{\emergencystretch}{3em} % prevent overfull lines

\providecommand{\tightlist}{%
  \setlength{\itemsep}{0pt}\setlength{\parskip}{0pt}}



 


\usepackage{booktabs}
\usepackage{longtable}
\usepackage{array}
\usepackage{multirow}
\usepackage{wrapfig}
\usepackage{float}
\usepackage{colortbl}
\usepackage{pdflscape}
\usepackage{tabu}
\usepackage{threeparttable}
\usepackage{threeparttablex}
\usepackage[normalem]{ulem}
\usepackage{makecell}
\usepackage{xcolor}
\usepackage{float}
\usepackage{placeins}
\usepackage[left]{lineno}
\linenumbers
\modulolinenumbers[1]
\KOMAoption{captions}{tableheading}
\makeatletter
\@ifpackageloaded{caption}{}{\usepackage{caption}}
\AtBeginDocument{%
\ifdefined\contentsname
  \renewcommand*\contentsname{Table of contents}
\else
  \newcommand\contentsname{Table of contents}
\fi
\ifdefined\listfigurename
  \renewcommand*\listfigurename{List of Figures}
\else
  \newcommand\listfigurename{List of Figures}
\fi
\ifdefined\listtablename
  \renewcommand*\listtablename{List of Tables}
\else
  \newcommand\listtablename{List of Tables}
\fi
\ifdefined\figurename
  \renewcommand*\figurename{Figure}
\else
  \newcommand\figurename{Figure}
\fi
\ifdefined\tablename
  \renewcommand*\tablename{Table}
\else
  \newcommand\tablename{Table}
\fi
}
\@ifpackageloaded{float}{}{\usepackage{float}}
\floatstyle{ruled}
\@ifundefined{c@chapter}{\newfloat{codelisting}{h}{lop}}{\newfloat{codelisting}{h}{lop}[chapter]}
\floatname{codelisting}{Listing}
\newcommand*\listoflistings{\listof{codelisting}{List of Listings}}
\makeatother
\makeatletter
\makeatother
\makeatletter
\@ifpackageloaded{caption}{}{\usepackage{caption}}
\@ifpackageloaded{subcaption}{}{\usepackage{subcaption}}
\makeatother
\usepackage{bookmark}
\IfFileExists{xurl.sty}{\usepackage{xurl}}{} % add URL line breaks if available
\urlstyle{same}
\hypersetup{
  pdftitle={d-CCSSBU L and S Charts and Poisson u Chart Code and Their Outputs for Equipment A and B},
  colorlinks=true,
  linkcolor={blue},
  filecolor={Maroon},
  citecolor={Blue},
  urlcolor={Blue},
  pdfcreator={LaTeX via pandoc}}


\title{d-CCSSBU L and S Charts and Poisson u Chart Code and Their
Outputs for Equipment A and B}
\author{}
\date{}
\begin{document}
\maketitle


\floatplacement{table}{H}
\floatplacement{figure}{H}

\section{Read Excel Data}\label{read-excel-data}

Each Excel Data contains three spreadsheet: DAR, CAR, and Mic. Each
spreadsheet should have variables \texttt{Event}, \texttt{Location},
\texttt{Result}, \texttt{USL}, \texttt{LOD}, \texttt{LOQ} (For Mic, it
does not have \texttt{LOD} and \texttt{LOQ}). You need to change the
following path to yours and download these data, code to your local
system. You need to have R/RStudio with necessary R packages installed.

You can also realize the algorithms in the manuscript directly in other
programming languages.

\section{d-CCSSBU-L Function}\label{d-ccssbu-l-function}

The function \texttt{d\_ccssbu\_l()} uses Excel data as above to
generate D-CCSSBU-L chart.

\section{\texorpdfstring{Selection of \(d\) for d-CCSSBU-L
Function}{Selection of d for d-CCSSBU-L Function}}\label{selection-of-d-for-d-ccssbu-l-function}

\begin{itemize}
\item
  \texttt{d\_ccssbu\_l()} constructs the d-CCSSBU-L chart for a
  specified value of \(d\).
\item
  \texttt{d\_ccssbu\_l\_d\_selection\_arl1()} applies
  \texttt{d\_ccssbu\_l()} across all candidate values \(d=1,\ldots,q\)
  using Excel input data and summarizes OOC counts and ARL\(_1\)
  performance for DAR, CAR, and Mic residue types.
\item
  \texttt{make\_d\_ccssbu\_l\_arl1\_wide\_tbl()} converts the output of
  \texttt{d\_ccssbu\_l\_d\_selection\_arl1()} into a wide-format summary
  table.
\item
  \texttt{render\_d\_ccssbu\_l\_arl1\_tbl()} renders the output table
  from \texttt{make\_d\_ccssbu\_l\_arl1\_wide\_tbl()} as a responsive
  tibble that automatically adjusts to page width.
\item
  \texttt{render\_d\_ccssbu\_l\_recommendation\_tbl()} generates the
  recommended \(d\) selection table for each residue type (DAR, CAR, and
  Mic).
\end{itemize}

\section{Application of the functions to Equipment A and
B.}\label{application-of-the-functions-to-equipment-a-and-b.}

\subsection{\texorpdfstring{d-CCSSBU-L Selection of \(d\) and ARL1
Results for Equipment
A}{d-CCSSBU-L Selection of d and ARL1 Results for Equipment A}}\label{d-ccssbu-l-selection-of-d-and-arl1-results-for-equipment-a}

Applying the \texttt{d\_ccssbu\_l\_d\_selection\_arl1()} function in
{[}\citeproc{ref-Yang2026dccssbu}{1}{]} to Equipment A data produces a
list of candidate values of \(d\) for each residue type. For a given
residue type (e.g., DAR), the output contains \(q\) rows, where \(q\) is
the number of ordered clusters identified.

For each candidate value of \(d\), the first \(d\) clusters are selected
to form the calibration dataset containing \texttt{nCal} cleaning
events. The calibration dataset is then used in the stratified smoothed
bootstrap procedure to estimate the chart UCL.

For the resulting UCL, the numbers of OOC events among ND events
(\texttt{OOC\_ND}), Stage 3A events (\texttt{OOC\_3A}), and Stage 3B
events (\texttt{OOC\_3B}) are reported. An
\texttt{OOC\_ND\ \textgreater{}\ 0} indicates that at least one ND event
is identified as OOC. Because ND events represent the best-quality
cleaning events, such a result suggests that the calibration dataset may
not be representative of the IC process. Similarly,
\texttt{OOC\_3A\ \textgreater{}\ 0} may indicate that the selected
calibration dataset does not adequately represent an IC Stage 3A
process.

For comparison with conventional control charts, the sampling-location
size associated with each cleaning event is recorded, and the median
value is used as the benchmark subgroup size. This benchmark subgroup
size is used to construct the conventional \(\bar{X}\) chart and \(S\)
chart and to calculate the corresponding benchmark \(\bar{X}\)-chart
ARL\(_1\) values reported in the table.

Finally, the d-CCSSBU-L chart ARL\(_1\) values are calculated for a
range of shift magnitudes. For the d-CCSSBU-L chart, the shift size is
defined as \(\delta_k = k \times \mathrm{UCL}/3\), where \(k\) denotes
the shift multiplier. The benchmark \(\bar{X}\)-chart ARL\(_1\) values
are provided for comparison with the proposed chart.

\begin{table}[!h]
\centering
\caption{\label{tab:display-d-selection-arl1-l-equipment-a-wide-v7}d-CCSSBU-L ARL1 simulation results for Equipment A. For each d, the d-CCSSBU-L chart was fitted using the location-specific calibration algorithm. ARL1 was estimated by shifting the bootstrap location statistics using delta k(d) equal to k times UCLboot(d) divided by 3.}
\centering
\resizebox{\ifdim\width>\linewidth\linewidth\else\width\fi}{!}{
\fontsize{7}{9}\selectfont
\begin{tabular}[t]{lrrlrrrrrrrrrrrrr}
\toprule
Residue & d & q & case\_label & nCal & OOC\_ND & OOC\_3A & OOC\_3B & benchmark\_n & Xbar\_chart\_ARL1\_k\_0p5 & UCL & ARL1\_k\_0p5 & ARL1\_k\_1 & ARL1\_k\_1p5 & ARL1\_k\_2 & ARL1\_k\_2p5 & ARL1\_k\_3\\
\midrule
CAR & 1 & 1 & Case 1 & 10 & 0 & 0 & 0 & 3 & 60.8923 & 4.1222 & 36.5631 & 36.5631 & 36.5631 & 35.9712 & 1.0000 & 1.0000\\
DAR & 1 & 3 & Case 2 & 3 & 0 & 0 & 0 & 6 & 26.3658 & 3.9132 & 6.8097 & 3.7064 & 2.2026 & 1.4708 & 1.0641 & 1.0000\\
DAR & 2 & 3 & Case 2 & 6 & 0 & 0 & 0 & 6 & 26.3658 & 4.5849 & 21.5285 & 7.2046 & 3.4764 & 2.0640 & 1.3702 & 1.0000\\
DAR & 3 & 3 & Case 2 & 10 & 0 & 0 & 0 & 6 & 26.3658 & 4.0912 & 50.7614 & 28.9855 & 9.1996 & 5.8428 & 2.5445 & 1.0002\\
Mic & 1 & 1 & Case 1 & 10 & 0 & 0 & 0 & 2 & 91.5158 & 2.0736 & 9.9751 & 9.9751 & 9.9751 & 9.9751 & 9.9751 & 1.2997\\
\bottomrule
\end{tabular}}
\end{table}

\FloatBarrier

The recommendation algorithm recommended CAR \(d\)=1, DAR \(d=2\) by
Case 1 in Step 8 although it totally has 3 clusters, and Mic \(d=1\).

\begin{table}[H]
\centering
\caption{\label{tab:display-d-selection-recommendation-l-equipment-a-v7}Recommended d values for the d-CCSSBU-L chart for Equipment A.}
\centering
\resizebox{\ifdim\width>\linewidth\linewidth\else\width\fi}{!}{
\fontsize{5}{7}\selectfont
\begin{tabular}[t]{lrrrlllllllll}
\toprule
Residue & recommended\_d & q & Xbar\_chart\_ARL1\_k\_0p5 & D1\_values & D2\_values & D3\_values & D4\_values & selected\_case & selection\_rule & D\_star\_values & selection\_path & recommendation\_notes\\
\midrule
CAR & 1 & 1 & 60.8923 & 1 & 1 & 1 &  & Case 1 & max(D1 intersect D2 intersect D3) & 1 & Case 1: max(D1 intersect D2 intersect D3) & \\
DAR & 2 & 3 & 26.3658 & 1, 2, 3 & 1, 2, 3 & 1, 2 & 3 & Case 1 & max(D1 intersect D2 intersect D3) & 1, 2 & Case 1: max(D1 intersect D2 intersect D3) & \\
Mic & 1 & 1 & 91.5158 & 1 & 1 & 1 &  & Case 1 & max(D1 intersect D2 intersect D3) & 1 & Case 1: max(D1 intersect D2 intersect D3) & \\
\bottomrule
\end{tabular}}
\end{table}

\FloatBarrier

\subsection{\texorpdfstring{d-CCSSBU-L Selection of \(d\) and ARL1
Results for Equipment
B}{d-CCSSBU-L Selection of d and ARL1 Results for Equipment B}}\label{d-ccssbu-l-selection-of-d-and-arl1-results-for-equipment-b}

The \texttt{d\_ccssbu\_l\_d\_selection\_arl1()} function was applied to
the Equipment B data using the same procedure described for Equipment A.

For CAR, the candidate solution with \(d=1\) produced OOC ND events
(\texttt{OOC\_ND\ =\ 3}). Although all OOC observations were ND, the
ordered-cluster algorithm identified two distinct groups of ND events,
corresponding to relatively smaller and larger ND values. When only the
first cluster was selected for calibration, the calibration dataset
became dominated by the smaller ND values, resulting in an excessively
small UCL. Consequently, several events containing larger ND values were
identified as OOC.

This result illustrates that a calibration dataset should be
representative of the IC process, even when all observations are ND.
Therefore, the candidate solution with \(d=1\) was not considered
suitable or the calibration dataset is not representative.

\begin{table}[!h]
\centering
\caption{\label{tab:display-d-selection-arl1-l-equipment-b-wide-v7}d-CCSSBU-L ARL1 simulation results for Equipment B. For each d, the d-CCSSBU-L chart was fitted using the location-specific calibration algorithm. ARL1 was estimated by shifting the bootstrap location statistics using delta k(d) equal to k times UCLboot(d) divided by 3.}
\centering
\resizebox{\ifdim\width>\linewidth\linewidth\else\width\fi}{!}{
\fontsize{7}{9}\selectfont
\begin{tabular}[t]{lrrlrrrrrrrrrrrrr}
\toprule
Residue & d & q & case\_label & nCal & OOC\_ND & OOC\_3A & OOC\_3B & benchmark\_n & Xbar\_chart\_ARL1\_k\_0p5 & UCL & ARL1\_k\_0p5 & ARL1\_k\_1 & ARL1\_k\_1p5 & ARL1\_k\_2 & ARL1\_k\_2p5 & ARL1\_k\_3\\
\midrule
CAR & 1 & 2 & Case 4 & 10 & 3 & 3 & 6 & 3 & 60.8923 & 2.1310 & 1.0000 & 1.0000 & 1.0000 & 1.0000 & 1.0000 & 1.0000\\
CAR & 2 & 2 & Case 4 & 20 & 0 & 0 & 0 & 3 & 60.8923 & 30.9387 & 500.0000 & 9.9701 & 6.3837 & 4.6083 & 3.5549 & 1.0000\\
DAR & 1 & 26 & Case 4 & 10 & 0 & 0 & 0 & 5 & 33.4221 & 19.5308 & 1.0088 & 1.0088 & 1.0088 & 1.0088 & 1.0088 & 1.0000\\
DAR & 2 & 26 & Case 5 & 13 & 0 & 0 & 0 & 5 & 33.4221 & 19.5454 & 1.2732 & 1.0051 & 1.0035 & 1.0032 & 1.0029 & 1.0000\\
DAR & 3 & 26 & Case 5 & 23 & 0 & 0 & 0 & 5 & 33.4221 & 23.4040 & 17.0648 & 3.4626 & 2.0206 & 1.3437 & 1.0264 & 1.0008\\
\addlinespace
DAR & 4 & 26 & Case 5 & 29 & 0 & 0 & 0 & 5 & 33.4221 & 20.7417 & 6.8399 & 3.9131 & 2.9291 & 1.9916 & 1.0162 & 1.0000\\
DAR & 5 & 26 & Case 5 & 39 & 0 & 0 & 0 & 5 & 33.4221 & 20.7252 & 19.0476 & 7.7942 & 4.6992 & 2.7941 & 1.1529 & 1.0004\\
DAR & 6 & 26 & Case 5 & 49 & 0 & 0 & 0 & 5 & 33.4221 & 20.4053 & 31.4961 & 12.5156 & 7.5472 & 4.0875 & 1.3605 & 1.0014\\
DAR & 7 & 26 & Case 5 & 59 & 0 & 0 & 0 & 5 & 33.4221 & 19.9548 & 44.0529 & 16.5289 & 11.4548 & 6.8236 & 1.7751 & 1.0003\\
DAR & 8 & 26 & Case 5 & 69 & 0 & 0 & 0 & 5 & 33.4221 & 19.6397 & 64.7249 & 22.3714 & 17.2117 & 10.2564 & 2.4393 & 1.0001\\
\addlinespace
DAR & 9 & 26 & Case 5 & 79 & 0 & 0 & 0 & 5 & 33.4221 & 19.5022 & 90.4977 & 30.3030 & 24.8756 & 14.8478 & 3.2321 & 1.0000\\
DAR & 10 & 26 & Case 5 & 89 & 12 & 9 & 0 & 5 & 33.4221 & 19.3354 & 133.3333 & 39.1389 & 34.5423 & 19.5122 & 4.2256 & 1.0001\\
DAR & 11 & 26 & Case 5 & 99 & 12 & 9 & 0 & 5 & 33.4221 & 19.2718 & 198.0198 & 51.9481 & 46.8384 & 27.1739 & 5.3262 & 1.0001\\
DAR & 12 & 26 & Case 5 & 109 & 12 & 9 & 0 & 5 & 33.4221 & 19.1111 & 312.5000 & 83.6820 & 76.3359 & 38.4615 & 8.2781 & 1.0009\\
DAR & 13 & 26 & Case 5 & 117 & 12 & 9 & 0 & 5 & 33.4221 & 18.9497 & 434.7826 & 97.5610 & 91.3242 & 42.2833 & 9.5283 & 1.0014\\
\addlinespace
DAR & 14 & 26 & Case 5 & 124 & 12 & 9 & 0 & 5 & 33.4221 & 18.9689 & 500.0000 & 124.2236 & 118.3432 & 51.8135 & 10.8342 & 1.0010\\
DAR & 15 & 26 & Case 5 & 126 & 12 & 9 & 0 & 5 & 33.4221 & 19.0971 & 434.7826 & 119.0476 & 109.2896 & 60.9756 & 11.7716 & 1.0012\\
DAR & 16 & 26 & Case 5 & 136 & 12 & 9 & 0 & 5 & 33.4221 & 18.6654 & 666.6667 & 136.9863 & 135.1351 & 59.1716 & 13.7741 & 1.0011\\
DAR & 17 & 26 & Case 5 & 140 & 12 & 9 & 0 & 5 & 33.4221 & 19.1061 & 500.0000 & 144.9275 & 139.8601 & 77.5194 & 15.3139 & 1.0008\\
DAR & 18 & 26 & Case 5 & 150 & 12 & 9 & 0 & 5 & 33.4221 & 18.4172 & 666.6667 & 188.6792 & 186.9159 & 66.2252 & 17.0358 & 1.0013\\
\addlinespace
DAR & 19 & 26 & Case 5 & 156 & 12 & 9 & 0 & 5 & 33.4221 & 14.2788 & 206.1856 & 206.1856 & 186.9159 & 37.3832 & 9.2166 & 1.0019\\
DAR & 20 & 26 & Case 5 & 162 & 12 & 9 & 0 & 5 & 33.4221 & 14.1476 & 210.5263 & 208.3333 & 194.1748 & 41.6667 & 10.6838 & 1.0023\\
DAR & 21 & 26 & Case 5 & 165 & 12 & 9 & 0 & 5 & 33.4221 & 16.2715 & 307.6923 & 202.0202 & 190.4762 & 44.4444 & 15.5280 & 1.0030\\
DAR & 22 & 26 & Case 5 & 175 & 15 & 10 & 0 & 5 & 33.4221 & 13.7396 & 289.8551 & 285.7143 & 263.1579 & 50.2513 & 10.9769 & 1.0036\\
DAR & 23 & 26 & Case 5 & 185 & 15 & 10 & 0 & 5 & 33.4221 & 13.9869 & 289.8551 & 289.8551 & 277.7778 & 53.0504 & 12.5000 & 1.0038\\
\addlinespace
DAR & 24 & 26 & Case 5 & 190 & 15 & 10 & 0 & 5 & 33.4221 & 13.9279 & 307.6923 & 307.6923 & 289.8551 & 61.1621 & 14.0647 & 1.0000\\
DAR & 25 & 26 & Case 5 & 199 & 15 & 10 & 0 & 5 & 33.4221 & 13.8748 & 317.4603 & 317.4603 & 307.6923 & 60.7903 & 14.7820 & 1.0000\\
DAR & 26 & 26 & Case 5 & 206 & 15 & 10 & 0 & 5 & 33.4221 & 13.6825 & 434.7826 & 425.5319 & 400.0000 & 76.9231 & 16.0256 & 1.0091\\
Mic & 1 & 1 & Case 4 & 10 & 0 & 0 & 0 & 3 & 60.8923 & 4.0806 & 28.6944 & 28.6944 & 28.6944 & 28.6944 & 28.6944 & 1.7109\\
\bottomrule
\end{tabular}}
\end{table}

\FloatBarrier

The following table summarizes the recommended values of \(d\) for
Equipment B. For CAR and Mic, only one admissible candidate remained
after applying the selection rules, resulting in recommended values of
\(d=2\) and \(d=1\), respectively. For DAR, multiple admissible
candidates were identified, and the selection procedure recommended
\(d=6\), corresponding to the smallest ARL\(_1\) among the admissible
candidates.

\begin{table}[H]
\centering
\caption{\label{tab:display-d-selection-recommendation-l-equipment-b-v7}Recommended d values for the d-CCSSBU-L chart for Equipment B.}
\centering
\resizebox{\ifdim\width>\linewidth\linewidth\else\width\fi}{!}{
\fontsize{5}{7}\selectfont
\begin{tabular}[t]{lrrrlllllllll}
\toprule
Residue & recommended\_d & q & Xbar\_chart\_ARL1\_k\_0p5 & D1\_values & D2\_values & D3\_values & D4\_values & selected\_case & selection\_rule & D\_star\_values & selection\_path & recommendation\_notes\\
\midrule
CAR & 2 & 2 & 60.8923 & 2 & 2 & 1 & 2 & Case 6 & min(D1 intersect D2 intersect D4) & 2 & Case 6: min(D1 intersect D2 intersect D4) & \\
DAR & 6 & 26 & 33.4221 & 1, 2, 3, 4, 5, 6, 7, 8, 9 & 1, 2, 3, 4, 5, 6, 7, 8, 9 & 1, 2, 3, 4, 5, 6 & 7, 8, 9, 10, 11, 12, 13, 14, 15, 16, 17, 18, 19, 20, 21, 22, 23, 24, 25, 26 & Case 1 & max(D1 intersect D2 intersect D3) & 1, 2, 3, 4, 5, 6 & Case 1: max(D1 intersect D2 intersect D3) & \\
Mic & 1 & 1 & 60.8923 & 1 & 1 & 1 &  & Case 1 & max(D1 intersect D2 intersect D3) & 1 & Case 1: max(D1 intersect D2 intersect D3) & \\
\bottomrule
\end{tabular}}
\end{table}

\FloatBarrier

\section{d-CCSSBU-S Function}\label{d-ccssbu-s-function}

The function \texttt{d\_ccssbu\_s()} uses Excel data as above to
generate D-CCSSBU-S chart.

\section{\texorpdfstring{Selection of \(d\) for d-CCSSBU-S
Function}{Selection of d for d-CCSSBU-S Function}}\label{selection-of-d-for-d-ccssbu-s-function}

\begin{itemize}
\item
  \texttt{d\_ccssbu\_s()} constructs the d-CCSSBU-L chart for a
  specified value of \(d\).
\item
  \texttt{d\_ccssbu\_s\_d\_selection\_arl1()} applies
  \texttt{d\_ccssbu\_s()} across all candidate values \(d=1,\ldots,q\)
  using Excel input data and summarizes OOC counts and ARL\(_1\)
  performance for DAR, CAR, and Mic residue types.
\item
  \texttt{make\_d\_ccssbu\_s\_arl1\_wide\_tbl()} converts the output of
  \texttt{d\_ccssbu\_s\_d\_selection\_arl1()} into a wide-format summary
  table.
\item
  \texttt{render\_d\_ccssbu\_s\_arl1\_tbl()} renders the output table
  from \texttt{make\_d\_ccssbu\_s\_arl1\_wide\_tbl()} as a responsive
  tibble that automatically adjusts to page width.
\item
  \texttt{render\_d\_ccssbu\_s\_recommendation\_tbl()} generates the
  recommended \(d\) selection table for each residue type (DAR, CAR, and
  Mic).
\end{itemize}

\section{\texorpdfstring{Application of the \(d\) Selection Functions to
Equipment A and
B.}{Application of the d Selection Functions to Equipment A and B.}}\label{application-of-the-d-selection-functions-to-equipment-a-and-b.}

\subsection{\texorpdfstring{d-CCSSBU-S Selection of \(d\) and ARL1
Results for Equipment
A}{d-CCSSBU-S Selection of d and ARL1 Results for Equipment A}}\label{d-ccssbu-s-selection-of-d-and-arl1-results-for-equipment-a}

Applying the \texttt{d\_ccssbu\_s\_d\_selection\_arl1()} function in
{[}\citeproc{ref-Yang2026dccssbu}{1}{]} to the Equipment A data produces
candidate values of \(d\) for each residue type. The d-selection
procedure is the same as that described for the d-CCSSBU-L chart, except
that the scale-specific calibration algorithm is used. For each
candidate value of \(d\), the corresponding OOC counts, UCL, benchmark
subgroup size, benchmark \(S\)-chart ARL\(_1\), and d-CCSSBU-S chart
ARL\(_1\) values are reported.

\begin{table}[!h]
\centering
\caption{\label{tab:display-d-selection-arl1-s-equipment-a-wide-v4}d-CCSSBU-S ARL1 simulation results for Equipment A. For each d, the d-CCSSBU-S chart was fitted using the scale-specific calibration algorithm. ARL1 was estimated by shifting the bootstrap scale statistics using delta k(d) equal to k times UCLboot(d) divided by 3.}
\centering
\resizebox{\ifdim\width>\linewidth\linewidth\else\width\fi}{!}{
\fontsize{7}{9}\selectfont
\begin{tabular}[t]{lrrlrrrrrrrrrrrrr}
\toprule
Residue & d & q & case\_label & nCal & OOC\_ND & OOC\_3A & OOC\_3B & benchmark\_n & S\_chart\_ARL1\_k\_0p5 & UCL & ARL1\_k\_0p5 & ARL1\_k\_1 & ARL1\_k\_1p5 & ARL1\_k\_2 & ARL1\_k\_2p5 & ARL1\_k\_3\\
\midrule
CAR & 1 & 1 & Case 1 & 10 & 0 & 0 & 0 & 3 & 98.2929 & 0.2011 & 206.1856 & 63.8978 & 22.8050 & 8.6393 & 2.8736 & 1.0000\\
DAR & 1 & 3 & Case 2 & 3 & 0 & 0 & 0 & 6 & 58.1871 & 4.6356 & 44.8430 & 18.9753 & 10.6496 & 3.3784 & 2.6438 & 1.0000\\
DAR & 2 & 3 & Case 2 & 7 & 0 & 0 & 0 & 6 & 58.1871 & 3.3661 & 540.5405 & 112.3596 & 29.8063 & 24.6305 & 10.4221 & 1.0003\\
DAR & 3 & 3 & Case 2 & 10 & 0 & 1 & 0 & 6 & 58.1871 & 3.1582 & 285.7143 & 87.3362 & 10.6952 & 6.5402 & 3.0492 & 1.0002\\
Mic & 1 & 1 & Case 1 & 10 & 0 & 0 & 0 & 2 & 132.1139 & 2.0516 & 9.9751 & 9.9751 & 9.9751 & 9.9751 & 9.9751 & 1.2997\\
\bottomrule
\end{tabular}}
\end{table}

\FloatBarrier

The following Table 6 summarizes the recommended values of \(d\) for the
d-CCSSBU-S chart for Equipment A. The recommended value was \(d=1\) for
CAR, DAR, and Mic. For DAR and Mic, the recommended candidate achieved a
smaller ARL\(_1(0.5)\) than the corresponding benchmark \(S\) chart
(Table 5).

For CAR, the recommended candidate did not achieve a smaller
ARL\(_1(0.5)\) than the benchmark \(S\) chart. This result is
attributable to the predominance of ND events in the calibration
dataset, which produced a very small bootstrap scale statistic and,
consequently, a very small UCL. Because the shift magnitude is defined
as \(\delta_k = k \times \mathrm{UCL}/3\), the resulting shift at
\(k=0.5\) is also very small. Consequently, the simulated OOC
distribution differs only slightly from the IC distribution, making
scale shifts difficult to detect and resulting in a large
ARL\(_1(0.5)\).

\begin{table}[!h]
\centering
\caption{\label{tab:display-d-selection-recommendation-s-equipment-a-v4}Recommended d values for the d-CCSSBU-S chart for Equipment A.}
\centering
\resizebox{\ifdim\width>\linewidth\linewidth\else\width\fi}{!}{
\fontsize{5}{7}\selectfont
\begin{tabular}[t]{lrrrlllllllll}
\toprule
Residue & recommended\_d & q & S\_chart\_ARL1\_k\_0p5 & D1\_values & D2\_values & D3\_values & D4\_values & selected\_case & selection\_rule & D\_star\_values & selection\_path & recommendation\_notes\\
\midrule
CAR & 1 & 1 & 98.2929 & 1 & 1 &  & 1 & Case 2 & min(D1 intersect D2 intersect D4) & 1 & Case 2: min(D1 intersect D2 intersect D4) & D3 empty: no candidate d satisfied ARL1(d, 0.5) less than or equal to the conventional S-chart benchmark.\\
DAR & 1 & 3 & 58.1871 & 1, 2, 3 & 1, 2 & 1 & 2, 3 & Case 1 & max(D1 intersect D2 intersect D3) & 1 & Case 1: max(D1 intersect D2 intersect D3) & \\
Mic & 1 & 1 & 132.1139 & 1 & 1 & 1 &  & Case 1 & max(D1 intersect D2 intersect D3) & 1 & Case 1: max(D1 intersect D2 intersect D3) & \\
\bottomrule
\end{tabular}}
\end{table}

\FloatBarrier

\subsection{\texorpdfstring{d-CCSSBU-S Selection of \(d\) and ARL1
Results for Equipment
B}{d-CCSSBU-S Selection of d and ARL1 Results for Equipment B}}\label{d-ccssbu-s-selection-of-d-and-arl1-results-for-equipment-b}

\begin{table}[!h]
\centering
\caption{\label{tab:display-d-selection-arl1-s-equipment-b-wide-v4}d-CCSSBU-S ARL1 simulation results for Equipment B. For each d, the d-CCSSBU-S chart was fitted using the scale-specific calibration algorithm. ARL1 was estimated by shifting the bootstrap scale statistics using delta k(d) equal to k times UCLboot(d) divided by 3.}
\centering
\resizebox{\ifdim\width>\linewidth\linewidth\else\width\fi}{!}{
\fontsize{7}{9}\selectfont
\begin{tabular}[t]{lrrlrrrrrrrrrrrrr}
\toprule
Residue & d & q & case\_label & nCal & OOC\_ND & OOC\_3A & OOC\_3B & benchmark\_n & S\_chart\_ARL1\_k\_0p5 & UCL & ARL1\_k\_0p5 & ARL1\_k\_1 & ARL1\_k\_1p5 & ARL1\_k\_2 & ARL1\_k\_2p5 & ARL1\_k\_3\\
\midrule
CAR & 1 & 2 & Case 4 & 10 & 0 & 0 & 0 & 3 & 98.2929 & 16.7963 & 96.1538 & 23.1750 & 7.5103 & 3.3881 & 1.7827 & 1.0092\\
CAR & 2 & 2 & Case 4 & 20 & 0 & 0 & 0 & 3 & 98.2929 & 14.2480 & 109.2896 & 22.3214 & 10.1626 & 6.3735 & 3.9331 & 1.0000\\
DAR & 1 & 26 & Case 4 & 10 & 0 & 1 & 2 & 5 & 67.2648 & 0.6525 & 181.8182 & 68.0272 & 24.3605 & 6.6028 & 1.7899 & 1.0000\\
DAR & 2 & 26 & Case 4 & 20 & 0 & 1 & 0 & 5 & 67.2648 & 19.1504 & 39.7614 & 39.7614 & 39.7614 & 29.4118 & 29.4118 & 1.0003\\
DAR & 3 & 26 & Case 5 & 23 & 0 & 1 & 0 & 5 & 67.2648 & 19.0995 & 63.2911 & 32.7332 & 30.4414 & 24.1255 & 6.8611 & 1.0004\\
\addlinespace
DAR & 4 & 26 & Case 5 & 33 & 0 & 1 & 0 & 5 & 67.2648 & 18.8003 & 111.7318 & 62.1118 & 51.8135 & 37.6648 & 12.1507 & 1.0021\\
DAR & 5 & 26 & Case 5 & 41 & 0 & 1 & 0 & 5 & 67.2648 & 18.6429 & 180.1802 & 90.9091 & 84.7458 & 63.6943 & 19.1571 & 1.0005\\
DAR & 6 & 26 & Case 5 & 51 & 0 & 1 & 0 & 5 & 67.2648 & 18.2699 & 322.5806 & 183.4862 & 141.8440 & 93.0233 & 30.5344 & 1.0018\\
DAR & 7 & 26 & Case 5 & 60 & 0 & 1 & 0 & 5 & 67.2648 & 17.5787 & 476.1905 & 208.3333 & 175.4386 & 101.5228 & 38.0228 & 1.0054\\
DAR & 8 & 26 & Case 5 & 63 & 0 & 1 & 0 & 5 & 67.2648 & 18.3464 & 312.5000 & 200.0000 & 168.0672 & 125.7862 & 42.6439 & 1.0053\\
\addlinespace
DAR & 9 & 26 & Case 5 & 73 & 0 & 1 & 0 & 5 & 67.2648 & 13.7839 & 298.5075 & 285.7143 & 232.5581 & 63.8978 & 39.7614 & 1.0107\\
DAR & 10 & 26 & Case 5 & 83 & 0 & 1 & 0 & 5 & 67.2648 & 13.4980 & 317.4603 & 298.5075 & 208.3333 & 16.2075 & 11.0988 & 1.0120\\
DAR & 11 & 26 & Case 5 & 89 & 0 & 1 & 0 & 5 & 67.2648 & 12.7163 & 588.2353 & 500.0000 & 134.2282 & 17.9372 & 7.0721 & 1.0128\\
DAR & 12 & 26 & Case 5 & 99 & 0 & 1 & 0 & 5 & 67.2648 & 12.6304 & 555.5556 & 434.7826 & 156.2500 & 22.6244 & 8.8928 & 1.0146\\
DAR & 13 & 26 & Case 5 & 101 & 0 & 1 & 0 & 5 & 67.2648 & 7.3237 & 96.1538 & 24.9688 & 18.5529 & 11.1545 & 4.2526 & 1.0115\\
\addlinespace
DAR & 14 & 26 & Case 5 & 111 & 0 & 1 & 0 & 5 & 67.2648 & 7.6154 & 259.7403 & 33.7268 & 25.0313 & 14.0252 & 5.0994 & 1.0122\\
DAR & 15 & 26 & Case 5 & 121 & 0 & 1 & 0 & 5 & 67.2648 & 10.5916 & 645.1613 & 434.7826 & 44.9438 & 25.1572 & 6.6979 & 1.0160\\
DAR & 16 & 26 & Case 5 & 128 & 0 & 1 & 0 & 5 & 67.2648 & 11.4061 & 666.6667 & 555.5556 & 66.8896 & 30.6748 & 8.5288 & 1.0219\\
DAR & 17 & 26 & Case 5 & 138 & 0 & 1 & 0 & 5 & 67.2648 & 7.2456 & 131.5789 & 47.9616 & 35.0877 & 17.1233 & 4.7382 & 1.0147\\
DAR & 18 & 26 & Case 5 & 148 & 0 & 1 & 0 & 5 & 67.2648 & 7.7705 & 400.0000 & 63.4921 & 47.3934 & 24.4798 & 5.9737 & 1.0000\\
\addlinespace
DAR & 19 & 26 & Case 5 & 153 & 0 & 1 & 0 & 5 & 67.2648 & 6.5651 & 81.9672 & 61.5385 & 39.7614 & 16.5700 & 4.9751 & 1.0113\\
DAR & 20 & 26 & Case 5 & 159 & 0 & 1 & 0 & 5 & 67.2648 & 8.4309 & 588.2353 & 84.0336 & 56.8182 & 31.1042 & 7.4267 & 1.0113\\
DAR & 21 & 26 & Case 5 & 169 & 0 & 1 & 0 & 5 & 67.2648 & 6.5147 & 81.9672 & 63.8978 & 44.4444 & 19.4742 & 5.8038 & 1.0090\\
DAR & 22 & 26 & Case 5 & 179 & 0 & 1 & 0 & 5 & 67.2648 & 6.6566 & 133.3333 & 89.2857 & 59.5238 & 25.5754 & 7.1048 & 1.0088\\
DAR & 23 & 26 & Case 5 & 183 & 0 & 1 & 0 & 5 & 67.2648 & 6.3784 & 99.0099 & 80.0000 & 54.4959 & 24.0674 & 6.7889 & 1.0083\\
\addlinespace
DAR & 24 & 26 & Case 5 & 189 & 0 & 1 & 0 & 5 & 67.2648 & 6.3870 & 122.6994 & 93.8967 & 59.3472 & 25.5428 & 7.1124 & 1.0000\\
DAR & 25 & 26 & Case 5 & 196 & 0 & 1 & 0 & 5 & 67.2648 & 6.4983 & 123.4568 & 94.3396 & 62.3053 & 29.6736 & 8.1800 & 1.0059\\
DAR & 26 & 26 & Case 5 & 206 & 0 & 1 & 0 & 5 & 67.2648 & 6.3374 & 142.8571 & 108.1081 & 73.8007 & 32.2581 & 7.6423 & 1.0000\\
Mic & 1 & 1 & Case 3 & 10 & 0 & 0 & 0 & 3 & 98.2929 & 3.9834 & 132.4503 & 132.4503 & 132.4503 & 132.4503 & 132.4503 & 1.8406\\
\bottomrule
\end{tabular}}
\end{table}

\FloatBarrier

\begin{table}[!h]
\centering
\caption{\label{tab:display-d-selection-recommendation-s-equipment-b-v4}Recommended d values for the d-CCSSBU-S chart for Equipment B.}
\centering
\resizebox{\ifdim\width>\linewidth\linewidth\else\width\fi}{!}{
\fontsize{5}{7}\selectfont
\begin{tabular}[t]{lrrrlllllllll}
\toprule
Residue & recommended\_d & q & S\_chart\_ARL1\_k\_0p5 & D1\_values & D2\_values & D3\_values & D4\_values & selected\_case & selection\_rule & D\_star\_values & selection\_path & recommendation\_notes\\
\midrule
CAR & 1 & 2 & 98.2929 & 1, 2 & 1, 2 & 1 & 2 & Case 1 & max(D1 intersect D2 intersect D3) & 1 & Case 1: max(D1 intersect D2 intersect D3) & \\
DAR & 2 & 26 & 67.2648 & 1, 2, 3, 4, 5, 6, 7, 8, 9, 10, 11, 12, 13, 14, 15, 16, 17, 18, 19, 20, 21, 22, 23, 24, 25, 26 &  & 2, 3 & 1, 4, 5, 6, 7, 8, 9, 10, 11, 12, 13, 14, 15, 16, 17, 18, 19, 20, 21, 22, 23, 24, 25, 26 & Case 10 & min(D1 intersect D3) & 2, 3 & Case 10: min(D1 intersect D3) & D2 empty: all candidate d values flagged at least one nominal Stage 3A event as OOC.\\
Mic & 1 & 1 & 98.2929 & 1 & 1 &  & 1 & Case 2 & min(D1 intersect D2 intersect D4) & 1 & Case 2: min(D1 intersect D2 intersect D4) & D3 empty: no candidate d satisfied ARL1(d, 0.5) less than or equal to the conventional S-chart benchmark.\\
\bottomrule
\end{tabular}}
\end{table}

\FloatBarrier

\section{d-CCSSBU-L And d-CCSSBU-S Charts of Equipment A and
B}\label{d-ccssbu-l-and-d-ccssbu-s-charts-of-equipment-a-and-b}

\subsection{d-CCSSBU Charts for Equipment A and B with Recommended d
Values}\label{d-ccssbu-charts-for-equipment-a-and-b-with-recommended-d-values}

The d-CCSSBU-L and d-CCSSBU-S charts for DAR using the recommended
values of \(d\) for Equipment A and Equipment B. Similar to the
conventional \(\bar{X}\)-chart and \(S\)-chart pair, the d-CCSSBU-L
chart monitors process location, whereas the d-CCSSBU-S chart monitors
within-event variability. Together, the two charts provide complementary
assessments of process stability.

Equipment A corresponds to Stage 3A data and Equipment B corresponds to
Stage 3B data. Previous analyses showed that conventional
\(\bar{X}\)-\(S\) charts and median-MAD charts may generate a large
number of OOC signals for heterogeneous cleaning validation data
{[}\citeproc{ref-Yang2026dccssbu}{1}{]}. In contrast, the d-CCSSBU
charts provide stable monitoring performance while accounting for data
heterogeneity.

The proposed charts are calibrated using d highest-risk ordered
clusters. The underlying principle is that if the highest-risk group
remains stable and in control, then lower-risk groups are also expected
to remain in control. This risk-based calibration strategy enables a
common control limit to be established for heterogeneous cleaning
validation data.

For Equipment B, no OOC signals were observed on the d-CCSSBU-L chart,
indicating stable process location performance. On the d-CCSSBU-S chart,
Event 10 was identified as OOC. Inspection of the corresponding DAR
observations (15.93, 0.11, 0.11, 15.93, and 32.95) shows substantial
within-event variability, resulting in a large MAD value. The d-CCSSBU-S
chart therefore identified this event as exhibiting unusually high
within-event variation.

\begin{figure}[H]
\centering
\begin{tabular}{cc}
\includegraphics[width=0.48\textwidth]{d_CCSSBU_L_DAR_Equipment_A.png} &
\includegraphics[width=0.48\textwidth]{d_CCSSBU_L_DAR_Equipment_B.png} \\

\includegraphics[width=0.48\textwidth]{d_CCSSBU_S_DAR_Equipment_A.png} &
\includegraphics[width=0.48\textwidth]{d_CCSSBU_s_DAR_Equipment_B.png} \\

\end{tabular}
\end{figure}

The following charts present the d-CCSSBU-L and d-CCSSBU-S charts for
CAR using the recommended values of \(d\) for Equipment A and Equipment
B. Similar to the conventional \(\bar{X}\)-\(S\) chart pair, the
d-CCSSBU-L chart monitors process location, whereas the d-CCSSBU-S chart
monitors within-event variability.

CAR data typically form a single cluster because most observations are
ND or close to ND. For Equipment B, two CAR groups corresponding to the
original cleaning process and the improved cleaning process were
intentionally included to demonstrate the capability of the proposed
framework to evaluate process changes. The d-CCSSBU charts show that the
improved cleaning process using the new cleaning agent remains
comparable with the original cleaning process and can be assessed using
a common control limit, despite the process change. Consistent with
expectations, most CAR observations in both groups were ND events.

For Equipment B, Event 19 was identified as an OOC signal by the
d-CCSSBU-L chart but not by the d-CCSSBU-S chart. The corresponding CAR
observations were 35.49, 31.99, and 24.80, yielding a median value of
31.99, which was substantially higher than those of the surrounding
events. However, the within-event variability remained relatively small,
as reflected by the modest MAD value. Consequently, the event was
detected as a location shift rather than a scale shift.

This example illustrates the complementary roles of the two charts. The
d-CCSSBU-L chart is sensitive to changes in process level, while the
d-CCSSBU-S chart is sensitive to changes in within-event variability.
Together, they provide a comprehensive assessment of cleaning process
stability.

\begin{figure}[H]
\centering
\begin{tabular}{cc}


\includegraphics[width=0.48\textwidth]{d_CCSSBU_L_CAR_Equipment_A.png} &
\includegraphics[width=0.48\textwidth]{d_CCSSBU_L_CAR_Equipment_B.png} \\

\includegraphics[width=0.48\textwidth]{d_CCSSBU_S_CAR_Equipment_A.png} &
\includegraphics[width=0.48\textwidth]{d_CCSSBU_s_CAR_Equipment_B.png} \\

\end{tabular}
\end{figure}

The following figure presents the d-CCSSBU-L and d-CCSSBU-S charts for
Mic using the recommended values of \(d\) for Equipment A and Equipment
B. Because microbial count data are homogeneous and contain
predominantly zero counts, both the location and scale statistics remain
near zero for most cleaning events.

For Equipment A, Event 7 contains microbial counts of 0 and 1.
Consequently, both the event location statistic and within-event
variability increase relative to the remaining events, producing a
visible spike on both the d-CCSSBU-L and d-CCSSBU-S charts. All other
events contain zero counts and therefore have location and scale
statistics close to zero.

For Equipment B, Event 42 contains microbial counts of 1 and 1. As a
result, the event location statistic is elevated relative to the
surrounding events, producing a spike on the d-CCSSBU-L chart. However,
because both observations are identical, the MAD is zero and no
corresponding increase is observed on the d-CCSSBU-S chart.

These examples illustrate the complementary roles of the two charts. The
d-CCSSBU-L chart detects shifts in event location, whereas the
d-CCSSBU-S chart detects changes in within-event variability. For
homogeneous microbial count data with predominantly zero counts, both
charts correctly reflect the underlying process characteristics.

\begin{figure}[H]
\centering
\begin{tabular}{cc}


\includegraphics[width=0.48\textwidth]{d_CCSSBU_L_Mic_Equipment_A.png} &
\includegraphics[width=0.48\textwidth]{d_CCSSBU_L_Mic_Equipment_B.png} \\

\includegraphics[width=0.48\textwidth]{d_CCSSBU_S_Mic_Equipment_A.png} &
\includegraphics[width=0.48\textwidth]{d_CCSSBU_s_Mic_Equipment_B.png} \\
\end{tabular}
\end{figure}

\section{Poisson u Chart Illustration for Mic
Data}\label{poisson-u-chart-illustration-for-mic-data}

\begin{itemize}
\tightlist
\item
  \texttt{prepare\_mic\_data()} prepares Mic data.
\item
  \texttt{poisson\_gof\_test()} tests if Mic data follows Poisson
  distribution.
\item
  \texttt{poisson\_overdispersion\_test()} tests if Mic data is
  over-dispersed.
\item
  \texttt{poisson\_mic\_u\_chart()} draw Poisson u-chart. It uses above
  three functions.
\end{itemize}

\subsubsection{Mic of Equipment A}\label{mic-of-equipment-a}

The following table shows that the Mic data for Equipment A satisfy the
Poisson assumptions, including the goodness-of-fit test and
overdispersion assessment. Therefore, a conventional Poisson \emph{u}
chart can be applied to this dataset.

Figure X compares the Poisson \emph{u} chart with the d-CCSSBU-L and
d-CCSSBU-S charts for the same data. The Poisson \emph{u} chart and the
d-CCSSBU-L chart produce nearly identical results, with UCL values of
2.0974 and 2.07, respectively. Both charts identify the same process
behavior and no OOC events are observed.

The close agreement is expected because the Mic data are homogeneous,
consist predominantly of zero counts, and satisfy the Poisson
distribution assumption. Under these conditions, the d-CCSSBU-L chart
effectively reproduces the performance of the conventional Poisson
\emph{u} chart.

For the same reason, the d-CCSSBU-L and d-CCSSBU-S charts are also
highly similar. Most cleaning events contain only zero counts, resulting
in both the location and scale statistics remaining close to zero
throughout the monitoring period. Consequently, the location and scale
charts exhibit nearly identical patterns.

\FloatBarrier

\begin{table}[H]
\centering
\caption{\label{tab:display-poisson-u-decision-summary-equipment-a-v3}Mic data Poisson assessment and chart selection for Equipment A.}
\centering
\resizebox{\ifdim\width>\linewidth\linewidth\else\width\fi}{!}{
\fontsize{8}{10}\selectfont
\begin{tabular}[t]{llrrlrllrrrrr}
\toprule
Equipment & Case & USL & GOF p & GOF pass & Dispersion & Overdispersion pass & Mic Poisson & CL & UCL min & UCL max & Stage 3A OOC & Stage 3B OOC\\
\midrule
Equipment A & Case 3 & 25 & 0.627 & TRUE & 1 & TRUE & TRUE & 0.2 & 2.0974 & 2.0974 & 0 & 0\\
\bottomrule
\end{tabular}}
\end{table}

\begin{center}
\includegraphics[width=0.95\linewidth,height=\textheight,keepaspectratio]{Poisson_U_Mic_Equipment_A.png}
\end{center}

\subsubsection{Mic of Equipment B}\label{mic-of-equipment-b}

The following table shows that the Mic data for Equipment B do not
satisfy the Poisson assumptions. The goodness-of-fit test failed
(\(p=0.0274\)), and the overdispersion assessment also failed with a
dispersion ratio of 1.6286. Consequently, the Mic data were not
considered Poisson, and a conventional Poisson \emph{u} chart was not
appropriate for this dataset.

This conclusion is also supported by the d-CCSSBU charts. Unlike
Equipment A, where the d-CCSSBU-L and d-CCSSBU-S charts were nearly
identical and closely matched the Poisson \emph{u} chart, the location
and scale charts for Equipment B exhibit different patterns. The
d-CCSSBU-L chart reflects differences in event-level microbial counts,
whereas the d-CCSSBU-S chart reflects within-event variability. The
visual differences between the two charts indicate that the underlying
process behavior is not consistent with a simple Poisson mechanism, for
which the mean and variance are intrinsically linked.

The chart results therefore agree with the formal statistical
assessment. Both the goodness-of-fit and overdispersion tests indicate
that the Poisson assumption is inadequate for the Equipment B Mic data,
while the differing behaviors of the d-CCSSBU-L and d-CCSSBU-S charts
provide a visual confirmation of this conclusion. Nevertheless, no OOC
events were identified for either Stage 3A or Stage 3B, indicating that
the microbial monitoring process remained stable throughout the study
period.

\begin{table}[!h]
\centering
\caption{\label{tab:display-poisson-u-decision-summary-equipment-b-v3}Mic data Poisson assessment and chart selection for Equipment B.}
\centering
\resizebox{\ifdim\width>\linewidth\linewidth\else\width\fi}{!}{
\fontsize{8}{10}\selectfont
\begin{tabular}[t]{llrrlrllrrrrr}
\toprule
Equipment & Case & USL & GOF p & GOF pass & Dispersion & Overdispersion pass & Mic Poisson & CL & UCL min & UCL max & Stage 3A OOC & Stage 3B OOC\\
\midrule
Equipment B & Case 3 & 25 & 0.0274 & FALSE & 1.6286 & FALSE & FALSE & 0.3333 & 1.488 & 2.3333 & 0 & 0\\
\bottomrule
\end{tabular}}
\end{table}

\protect\phantomsection\label{refs}
\begin{CSLReferences}{0}{1}
\bibitem[\citeproctext]{ref-Yang2026dccssbu}
\CSLLeftMargin{1. }%
\CSLRightInline{Yang CX. D-CCSSBU-l-or-s control chart for cleaning
validation {[}Internet{]}. GitHub repository; 2026. Available from:
\url{https://github.com/ChandlerXiandeYang/d-CCSSBU-L-or-S-Control-Chart}}

\end{CSLReferences}




\end{document}
